\chapter{Navigation in the Mass Framework}
\label{sec:Navigation}

Navigating a massive number of agents through a complex 3D environment is traditionally one of the most CPU-intensive tasks in game development. In Unreal Engine's standard Object-Oriented paradigm, navigation relies heavily on \texttt{AController} and \texttt{UPathFollowingComponent}. While robust, this approach requires constant pointer-chasing and individual tick updates, creating a hard ceiling on the number of autonomous agents the engine can simulate simultaneously.

Within the Mass framework, navigation is completely reimagined to fit the Entity Component System (ECS) architecture. This chapter explores the navigation paradigms available in Mass, detailing the data-oriented implementations of Navigation Mesh (NavMesh) and ZoneGraph, demonstrating how to build hybrid navigation systems, and examining the low-level steering and avoidance processors.

\section{Overview}

Mass does not entirely reinvent the mathematical wheel of pathfinding. Instead, it optimizes how path requests are batched, stored, and executed. The MassNavigation and MassAI plugins provide the structural logic for moving entities, decoupling the intent to move from the physical execution of movement.

Regardless of the chosen pathfinding method, the ultimate goal of the Mass navigation pipeline is to populate and continuously update an entity's \newline\texttt{FMassMoveTargetFragment}. This universal fragment acts as the movement intent, storing the immediate destination vector, desired speed, and distance to the goal. Processors downstream read this fragment without needing to know whether the target was generated by a complex NavMesh query or a simple spline calculation.

\section{ZoneGraph}

While NavMesh is highly versatile for open-world traversal, computing routes across thousands of interconnected polygons is computationally expensive. For simulations like the City Sample, where thousands of pedestrians and vehicles follow strict traffic laws, Epic Games introduced ZoneGraph.

%immagine ZoneGraph

ZoneGraph is a spline-based routing system completely detached from the NavMesh. Instead of a walkable surface area, ZoneGraph defines explicit lanes, intersections, and connections. It operates on a conceptual level similar to a rail network, built upon four foundational structural elements that we will now analyse.

\subsection{Zones}

Zones are the high-level containers that define the overall trajectory of a route. In the Unreal Editor, a Zone is the primary authoring unit, typically represented by a spline curve. However, it does not exist as an Actor at runtime. During the environment baking process, these splines are compiled into flattened data structures (such as \texttt{FZoneGraphStorage}). A Zone does not dictate specific agent movement. Instead, it provides the geometric foundation for grouping multiple parallel tracks. This grouping helps Mass spatial partitioning systems efficiently manage entity Level of Detail (LOD) logic based on the agent's current sector.

\subsection{Lanes}

Lanes are the specific, functional tracks within a Zone that dictate the exact path an entity follows. A single Zone can contain multiple parallel lanes, each with distinct semantic profiles, tags, and traversal rules (e.g., separating a forward vehicle lane from a pedestrian sidewalk). At runtime, an entity's position on a lane is tracked using lightweight data fragments, specifically the \texttt{FMassZoneGraphLaneLocationFragment}. This fragment stores the current lane index and the distance along its curve. By reducing movement to a 1D distance value along a track, Mass processors can advance the entity's position through simple scalar addition, bypassing the need for continuous 3D surface projections or complex path smoothing.

\subsection{Intersections} 

They act as connective nodes that link individual lanes together to form a cohesive network. Rather than relying on real-time spatial queries to discover the next path, intersections provide entities with a predefined set of linked lane indices. They act as logical checkpoints where the Mass framework can handle branching routing (like turning at a crossroad) and StateTree tasks can evaluate conditional flow (like yielding at a red light) through highly efficient array lookups.

\subsection{Connections}

Connections are the explicit data links that define how individual lanes relate to one another. While intersections serve as the major routing hubs, connections represent the directed, point-to-point transitions, such as moving forward into a new spline segment or shifting laterally to an adjacent lane. Within the ECS architecture, these relationships are stored as simple arrays of linked lane indices (e.g., \texttt{FZoneGraphLinkedLane}). By pre-calculating these valid transitions during the environment baking process, the framework allows entities to execute complex maneuvers, like overtaking a slower vehicle or merging into traffic, through instantaneous memory lookups. This entirely eliminates the need for expensive physics sweeps or dynamic overlap checks when an agent needs to change its current track.\newline

A critical advantage of this entire system lies in its Data-Oriented Design. Unlike standard Unreal Engine splines that are typically constructed from multiple object-oriented Actors, ZoneGraph data, including all lanes and intersection links, is baked into flat, contiguous memory arrays. This structural alignment allows the navigation data to be queried in constant time, making it exceptionally fast for ECS processors to read during parallel execution. Ultimately, because the path is explicitly constrained by the lane's mathematical geometry and intersection rules, the complex burden of spatial pathfinding is reduced to a highly optimized, one-dimensional graph-traversal problem.

\section{NavMesh Navigation}

%RIPARTI DA QUI, L'HA SCRITTO GEMINI
For open-world environments, unstructured terrain, and dynamic spatial queries, Mass provides a data-oriented wrapper around the standard Unreal Engine Recast Navigation Mesh (dtNavMesh). Because querying the underlying NavMesh octree is an asynchronous, object-oriented operation, it cannot be safely executed within a high-speed, parallelized processor loop. Mass solves this by strictly separating the path request from the continuous path execution.

The NavMesh Trait
To utilize NavMesh routing, an entity's archetype must be configured with a specific movement trait via the editor, such as UMassNavMeshMovementTrait. This trait acts as the structural blueprint, ensuring the entity is initialized with the required memory containers. Most notably, it appends the FMassNavMeshShortPathFragment, a lightweight data structure designed to store the array of spatial waypoints and polygon references (dtPolyRef) that make up the navigation corridor.

The StateTree Task
The initiation of a pathfinding request is strictly driven by the AI's StateTree. When an entity needs to move to a dynamic location, the tree executes a task like FMassFindNavMeshPathTask. This task serves as the bridge between the ECS architecture and the traditional engine subsystems:

The Query: The task interfaces directly with UNavigationSystemV1, queuing an asynchronous path request to the destination.

The Wait State: While the standard navigation system computes the route, the task holds the StateTree execution in an "In Progress" state, preventing the entity from executing subsequent behavioral logic prematurely.

Data Translation: Upon receiving the successful callback, the task destructures the heavy, object-based FNavPathSharedPtr. It extracts the essential coordinate data and writes it directly into the entity's FMassNavMeshShortPathFragment, marking the task as successful.

The Path Following Processor
Once the path data is securely written to the chunk memory, the continuous logic of traversing the corridor is handed off to the UMassNavMeshPathFollowProcessor. This processor does not apply physical velocity directly; instead, it acts as a geometric interpolator.

During its parallel execution phase, the processor reads the FMassNavMeshShortPathFragment and calculates the entity's true progress by projecting its physical location onto the closest path segment. Using the entity's desired speed and the frame's delta time, it calculates a "look-ahead" target along the path's curve. It then pushes this exact spatial coordinate and tangent vector into the FMassMoveTargetFragment.

By continuously dragging this target point just ahead of the entity, the processor decouples the complexity of the NavMesh corridor from the actual physical locomotion. Furthermore, when the processor detects that the entity has reached the final waypoint within a predefined tolerance, it integrates with the UMassSignalSubsystem to broadcast a completion signal (e.g., UE::Mass::Signals::FollowPointPathDone). This signal wakes the sleeping StateTree task, prompting the AI to immediately evaluate its next behavioral state without requiring expensive, per-frame polling.